% GENERATED FILE. Edit canonical lessons, then run npm run generate:books.
% canonical-source-hash: fnv1a64:44a434e95362cbeb
% canonical-lessons: IT-C275-accadere, IT-C275-capitare, IT-C275-verificarsi, IT-C275-avvenire, IT-C275-accaduto

\chapter{Cosa è accaduto}
\label{ch:it-cosa-e-accaduto}
\hlchaptermodality{275}

\begin{chapteropening}
\textbf{By the end of this chapter:} \emph{I can say that something happened, came about or took place, and talk about what happened.}

Complete the last lesson of chapter 275: I can say that something happened, came about or took place, and talk about what happened.
\end{chapteropening}

\section[accadere]{accadere — to happen}

\label{lesson:IT-C275-accadere}



\begin{quote}
Before the new one: say the Italian for at a certain point, then the Italian for all of a sudden.
\end{quote}

\subsection*{You'll want to know: accadere}
\textbf{accadere} — \textquotedblleft{}to happen\textquotedblright{}.

\begin{quote}
È accaduto tutto in pochi secondi.
\end{quote}

\begin{quote}
\emph{It all happened in a few seconds.}
\end{quote}

Takes \emph{essere}: \emph{è accaduto}. A little more formal than \emph{succedere}.

\begin{cousinweb}
\emph{accadere} is usually traced to Vulgar Latin \emph{accadĕre}, a remake of Latin \emph{accidere} \textquotedblleft{}to fall upon, to happen\textquotedblright{}, from \emph{ad} and \emph{cadere} \textquotedblleft{}to fall\textquotedblright{} (Lewis \& Short).
\end{cousinweb}

\begin{grammarlens}[title={What you've built}]
One more word for telling, explaining and describing at B1.
\end{grammarlens}

\subsection*{Your turn}
\begin{itemize}
\raggedright
  \item \emph{Say it:} \emph{accadere}
  \item \emph{Say it:} \emph{accadere}, once more
  \item \emph{Say it:} a sentence of your own with \emph{accadere}
  \item \emph{Recall:} say the Italian for at a certain point, then the Italian for all of a sudden, then say \emph{accadere} again
\end{itemize}

\begin{tcolorbox}[breakable,colback=teal!4,colframe=teal!35!black,title={Before you move on}]
What does accadere mean? (\textquotedblleft{}To happen\textquotedblright{}.) Say it once more.
\end{tcolorbox}
\section[capitare]{capitare — to happen to someone; to turn up}

\label{lesson:IT-C275-capitare}



\begin{quote}
Before the new one: say the Italian for all of a sudden, then the Italian for to happen.
\end{quote}

\subsection*{You'll want to know: capitare}
\textbf{capitare} — \textquotedblleft{}to happen to someone; to turn up\textquotedblright{}.

\begin{quote}
Mi è capitato di incontrarlo in treno.
\end{quote}

\begin{quote}
\emph{I happened to meet him on the train.}
\end{quote}

\emph{Capita!} — these things happen. Takes \emph{essere}.

\begin{cousinweb}
\emph{capitare} is inherited from Vulgar Latin \emph{capitāre} \textquotedblleft{}to come to a head, to arrive\textquotedblright{}, built on Latin \emph{caput} \textquotedblleft{}head\textquotedblright{} (EtymDB).
\end{cousinweb}

\begin{grammarlens}[title={What you've built}]
One more word for telling, explaining and describing at B1.
\end{grammarlens}

\subsection*{Your turn}
\begin{itemize}
\raggedright
  \item \emph{Say it:} \emph{capitare}
  \item \emph{Say it:} \emph{capitare}, once more
  \item \emph{Say it:} a sentence of your own with \emph{capitare}
  \item \emph{Recall:} say the Italian for all of a sudden, then the Italian for to happen, then say \emph{capitare} again
\end{itemize}

\begin{tcolorbox}[breakable,colback=teal!4,colframe=teal!35!black,title={Before you move on}]
What does capitare mean? (\textquotedblleft{}To happen to someone; to turn up\textquotedblright{}.) Say it once more.
\end{tcolorbox}
\section[verificarsi]{verificarsi — to occur, to take place (formal)}

\label{lesson:IT-C275-verificarsi}



\begin{quote}
Before the new one: say the Italian for to happen, then the Italian for to happen to someone.
\end{quote}

\subsection*{You'll want to know: verificarsi}
\textbf{verificarsi} — \textquotedblleft{}to occur, to take place (formal)\textquotedblright{}.

\begin{quote}
Ieri sera si è verificato un incidente in centro.
\end{quote}

\begin{quote}
\emph{Last night an accident occurred in the centre.}
\end{quote}

A news word. Without \emph{-si}, \emph{verificare} means ‘to check’.

\begin{cousinweb}
\emph{verificarsi} comes from \emph{verificare} \textquotedblleft{}to check, to show true\textquotedblright{}, usually traced to Medieval Latin \emph{vērificāre}, built on \emph{vērus} \textquotedblleft{}true\textquotedblright{} and \emph{facere} \textquotedblleft{}to make\textquotedblright{} (Lewis \& Short).
\end{cousinweb}

\begin{grammarlens}[title={What you've built}]
One more word for telling, explaining and describing at B1.
\end{grammarlens}

\subsection*{Your turn}
\begin{itemize}
\raggedright
  \item \emph{Say it:} \emph{verificarsi}
  \item \emph{Say it:} \emph{verificarsi}, once more
  \item \emph{Say it:} a sentence of your own with \emph{verificarsi}
  \item \emph{Recall:} say the Italian for to happen, then the Italian for to happen to someone, then say \emph{verificarsi} again
\end{itemize}

\begin{tcolorbox}[breakable,colback=teal!4,colframe=teal!35!black,title={Before you move on}]
What does verificarsi mean? (\textquotedblleft{}To occur, to take place (formal)\textquotedblright{}.) Say it once more.
\end{tcolorbox}
\section[avvenire]{avvenire — to take place, to come about}

\label{lesson:IT-C275-avvenire}



\begin{quote}
Before the new one: say the Italian for to happen to someone, then the Italian for to occur.
\end{quote}

\subsection*{You'll want to know: avvenire}
\textbf{avvenire} — \textquotedblleft{}to take place, to come about\textquotedblright{}.

\begin{quote}
Il fatto è avvenuto di notte.
\end{quote}

\begin{quote}
\emph{The incident took place at night.}
\end{quote}

Conjugated like \emph{venire}: \emph{è avvenuto}. As a noun, \emph{l'avvenire} is the future.

\begin{cousinweb}
\emph{avvenire} is inherited from Latin \emph{advenīre} \textquotedblleft{}to come to, to arrive\textquotedblright{} (EtymDB), from \emph{ad} and \emph{venīre} \textquotedblleft{}to come\textquotedblright{}.
\end{cousinweb}

\begin{grammarlens}[title={What you've built}]
One more word for telling, explaining and describing at B1.
\end{grammarlens}

\subsection*{Your turn}
\begin{itemize}
\raggedright
  \item \emph{Say it:} \emph{avvenire}
  \item \emph{Say it:} \emph{avvenire}, once more
  \item \emph{Say it:} a sentence of your own with \emph{avvenire}
  \item \emph{Recall:} say the Italian for to happen to someone, then the Italian for to occur, then say \emph{avvenire} again
\end{itemize}

\begin{tcolorbox}[breakable,colback=teal!4,colframe=teal!35!black,title={Before you move on}]
What does avvenire mean? (\textquotedblleft{}To take place, to come about\textquotedblright{}.) Say it once more.
\end{tcolorbox}
\section[l'accaduto]{l'accaduto — what happened, the incident}

\label{lesson:IT-C275-accaduto}



\begin{quote}
Before the new one: say the Italian for to occur, then the Italian for to take place.
\end{quote}

\subsection*{You'll want to know: l'accaduto}
\textbf{l'accaduto} — \textquotedblleft{}what happened, the incident\textquotedblright{}.

\begin{quote}
Mi ha raccontato l'accaduto.
\end{quote}

\begin{quote}
\emph{He told me what had happened.}
\end{quote}

The past participle of \emph{accadere} used as a noun; common on the news.

\begin{cousinweb}
\emph{l'accaduto} is the past participle of \emph{accadere} used as a noun, \textquotedblleft{}the thing that happened\textquotedblright{}; the verb goes back to Latin \emph{cadere} \textquotedblleft{}to fall\textquotedblright{} (Lewis \& Short).
\end{cousinweb}

\begin{grammarlens}[title={What you've built}]
One more word for telling, explaining and describing at B1.
\end{grammarlens}

\subsection*{Your turn}
\begin{itemize}
\raggedright
  \item \emph{Say it:} \emph{l'accaduto}
  \item \emph{Say it:} \emph{l'accaduto}, once more
  \item \emph{Say it:} a sentence of your own with \emph{l'accaduto}
  \item \emph{Recall:} say the Italian for to occur, then the Italian for to take place, then say \emph{l'accaduto} again
\end{itemize}

\begin{tcolorbox}[breakable,colback=teal!4,colframe=teal!35!black,title={Before you move on}]
What does l'accaduto mean? (\textquotedblleft{}What happened, the incident\textquotedblright{}.) Say it once more.
\end{tcolorbox}
